\section{计算几何}

\subsection{二维几何基础}
	\subsubsection{声明与宏}
		\inputminted{cpp}{sections/ComputationalGeometry/assets/Geometry/8.1-declarations.cpp}
	\subsubsection{点与向量}
		\inputminted{cpp}{sections/ComputationalGeometry/assets/Geometry/8.2-point-vector.cpp}
	\subsubsection{线}
		\inputminted{cpp}{sections/ComputationalGeometry/assets/Geometry/8.3-line.cpp}
	\subsubsection{三角形}
		\inputminted{cpp}{sections/ComputationalGeometry/assets/Geometry/三角形.cpp}
	\subsubsection{三角形相关公式：三角形心、海伦公式、欧拉定理、内接球、外接圆}
		\input{sections/ComputationalGeometry/assets/Geometry/Formula(Saga).tex}

\subsection{多边形}
	\subsubsection{多边形基础}
		\inputminted{cpp}{sections/ComputationalGeometry/assets/Geometry/多边形基础.cpp}
	\subsubsection{简单多边形内最长线段}
		\inputminted{cpp}{sections/ComputationalGeometry/assets/Geometry/AirportConstruction.cpp}
	\subsubsection{$O(n ^ 2)$ Ear Clipping 三角剖分}
		\inputminted{cpp}{sections/ComputationalGeometry/assets/Geometry/Triangulation.cpp}

\subsection{圆}
	\subsubsection{圆基础}
		\inputminted{cpp}{sections/ComputationalGeometry/assets/Geometry/8.4-circle.cpp}
	\subsubsection{圆反演与阿波罗尼茨圆}
		\begin{small}
			\input{sections/ComputationalGeometry/assets/Geometry/阿波罗尼茨圆.tex}
		\end{small}
		\inputminted{cpp}{sections/ComputationalGeometry/assets/yzh/circle_inversion.cpp}
	\subsubsection{圆并}
		\inputminted{cpp}{sections/ComputationalGeometry/assets/Geometry/圆并.cpp}
	\subsubsection{多边形与圆交}
		\inputminted{cpp}{sections/ComputationalGeometry/assets/Geometry/多边形和圆的交.cpp}
	% \subsubsection{圆幂、圆反演、根轴}
	% 	\input{sections/ComputationalGeometry/assets/Geometry/圆幂.tex}
	\subsubsection{圆上整点}
		\inputminted{cpp}{sections/ComputationalGeometry/assets/Geometry/圆上整点.cpp}

\subsection{凸几何}
	\subsubsection{凸包}
		\inputminted{cpp}{sections/ComputationalGeometry/assets/Geometry/凸包.cpp}
	\subsubsection{凸包询问：凸包内、切点、交点、最近点}
		\inputminted{cpp}{sections/ComputationalGeometry/assets/Geometry/convex_findmax.cpp}
	\subsubsection{单插入动态凸包}
		\inputminted{cpp}{sections/ComputationalGeometry/assets/Geometry/dynamic_convex_hull.cpp}
		% \inputminted{cpp}{sections/ComputationalGeometry/assets/Geometry/logconvexhull.cpp}
	\subsubsection{旋转卡壳}
		\inputminted{cpp}{sections/ComputationalGeometry/assets/Geometry/旋转卡壳.cpp}
	\subsubsection{闵可夫斯基和}
		\inputminted{cpp}{sections/ComputationalGeometry/assets/Geometry/闵可夫斯基和.cpp}
	\subsubsection{直线半平面交}
		\inputminted{cpp}{sections/ComputationalGeometry/assets/Geometry/半平面交.cpp}
	\subsubsection{整数半平面交}
		\inputminted{cpp}{sections/ComputationalGeometry/assets/Geometry/integral_hpi.cpp}
	% \subsubsection{直线与凸包交点}
	% 	\inputminted{cpp}{sections/ComputationalGeometry/assets/Geometry/直线与凸包交点.cpp}
	% \subsubsection{点到凸包切线}
	% 	\inputminted{cpp}{sections/ComputationalGeometry/assets/Geometry/点到凸包切线.cpp}

\subsection{点集算法}
	\subsubsection{最近点对}
		\inputminted{cpp}{sections/ComputationalGeometry/assets/Geometry/最近点对.cpp}
	\subsubsection{最小覆盖球}
		\inputminted{cpp}{sections/ComputationalGeometry/assets/Geometry/最小覆盖球.cpp}
	\subsubsection{Delaunay 三角剖分}
		\inputminted{cpp}{sections/ComputationalGeometry/assets/Geometry/DelaunayTriangulation.cpp}
	% \subsubsection{长方体表面两点最短距离}
	% 	\inputminted{cpp}{sections/ComputationalGeometry/assets/Geometry/长方体表面两点最短距离.cpp}

\subsection{球面与三维几何}
	\subsubsection{球面基础与经纬度球面距离}
		\begin{small}
			\input{sections/ComputationalGeometry/assets/Geometry/球面.tex}
		\end{small}
		\inputminted{cpp}{sections/ComputationalGeometry/assets/Geometry/经纬度求球面最短距离.cpp}
	\subsubsection{三维几何基础操作}
		\inputminted{cpp}{sections/ComputationalGeometry/assets/Geometry/三维几何.cpp}
	\subsubsection{三维凸包}
		\inputminted{cpp}{sections/ComputationalGeometry/assets/Geometry/三维凸包.cpp}

\subsection{数值工具}
	\subsubsection{黄金三分}
		\inputminted[highlightlines={7}]{cpp}{sections/ComputationalGeometry/assets/yzh/golden_ternary.cpp}
	\subsubsection{自适应 Simpson}
		\inputminted{cpp}{sections/ComputationalGeometry/assets/Math/Simpson.cpp}
